\documentclass[10pt,a4paper]{amsart}
\usepackage[margin=19mm]{geometry}
\usepackage{amsmath,amssymb,mathtools}
\usepackage[scr=rsfs,cal=euler]{mathalfa}
\usepackage{xfrac}
\usepackage[unicode,hidelinks,bookmarks=false]{hyperref}
\newcommand{\ie}{i.e.\ }
\newcommand{\cf}{cf.\ }
\newcommand{\resp}{respectively\ }
\newcommand{\Subsection}{\subsection}
\providecommand{\nobreakdash}{\nobreak\hbox{-}}
\setlength{\emergencystretch}{2em}
\multlinegap0pt
\usepackage{bm}
\usepackage{bbm}
\usepackage{dsfont}
\usepackage{xspace}
\usepackage{cancel}
\usepackage{mathtools}
\usepackage{amsthm}
\makeatletter
\@ifundefined{lemma}{\newtheorem{lemma}{Lemma}}{}
\@ifundefined{prop}{\newtheorem{prop}{Proposition}}{}
\makeatother
\title[A Non-Autonomous Model for Parabolic Implosion]{A Non-Autonomous Model for Parabolic Implosion}
\author{Katelynn Huneycutt, Samantha Sandberg-Clark, Liz Vivas}
\date{Source: arXiv:2601.16311v1 (CC BY 4.0)}
\begin{document}
\maketitle

\noindent\emph{Source and licence.} A Non-Autonomous Model for Parabolic Implosion. Version arXiv:2601.16311v1 (CC BY 4.0), distributed under the Creative Commons Attribution 4.0 International licence, as recorded in the arXiv licence metadata for this exact version and shown on the abstract page https://arxiv.org/abs/2601.16311v1. Full rights notice: licenses/arxiv-2601.16311-CC-BY-4.0.txt.\par\smallskip

\noindent\textit{Editorial scope.} Counted unit: the Prop whose source label is \texttt{prop:randombigO} in the article (source0.tex lines 1311--1321, source label \texttt{prop:randombigO}), delivered together with the article's own complete original proof of it (source0.tex lines 1323--1402, one \texttt{proof} environment of 80 lines). No mathematics is written, repaired or re-derived: statement and proof are the article's own text line for line. Required context retained uncounted: lem:martingale (lines 1253--1260); lem:applyAzuma (lines 1284--1298). Every definition, display and label of those lines is retained unchanged. Omitted from this package and not redistributed: 1--1252, 1261--1283, 1299--1310, 1322--1322, 1403--1466. Nothing inside a retained excerpt is omitted or altered: every retained body is delivered exactly as the article's lines are written, so no omission receipt of any kind accompanies this package. Citations: the retained excerpts carry no citation command at all, so no bibliography entry of the article is transcribed and no reference is redistributed. Reference adaptation: none was needed and none was made; every reference inside the delivered unit resolves inside it, so no unresolved reference or citation is produced. Printed slips kept literal: no printed word, symbol, hypothesis or number of the counted statement or of its proof was altered. Layout and class: the article's own class is \texttt{amsart} with packages including geometry, mathtools, amsmath, amsthm, graphicx, caption, xcolor, hyperref, float, comment, mathrsfs, enumitem, dsfont; the wrapper is the lane's pinned \texttt{amsart} editorial wrapper at 10pt on A4, which restates the theorem-like environments the retained text needs and adds the article's own macro definitions that the retained text needs (none), with extra wrapper packages (bm, bbm, dsfont, xspace, cancel, mathtools). The delivered numbering is the unit's own. Assets: no retained span contains a figure environment, a graphics-inclusion command, a PDF-page inclusion, a table or a TikZ picture; no image file is copied into this package, the package declares no asset dependency and no image file is redistributed with it. Licence: version arXiv:2601.16311v1 of this article is distributed under the Creative Commons Attribution 4.0 International licence, as recorded in the arXiv licence metadata for this exact version and shown on the abstract page https://arxiv.org/abs/2601.16311v1. A full rights notice is delivered at licenses/arxiv-2601.16311-CC-BY-4.0.txt. Characters: every retained excerpt is pure ASCII, so the delivered engine, XeLaTeX with its default Latin Modern text font, reports no missing character.
\begin{lemma}\label{lem:martingale}
Let
\[
\delta_n:=\sum_{k=0}^{n-1} d_k\,q_k\,e^{ik\theta}, \qquad n\ge1.
\]
Then $\{\delta_n\}_{n\ge1}$ is a martingale with respect to the filtration
$\mathcal{F}_n=\sigma(\xi_0,\dots,\xi_{n-1})$.
\end{lemma}

\begin{lemma}\label{lem:applyAzuma}
Assume $|\xi_k|\le M$ almost surely and $|q_k|\le N$ with probability one, for $k\le n$.  
Then for any $\lambda_n>0$,
\[
\mathbb{P}(|\delta_n|\ge\lambda_n)
\le
\exp\!\left(-\frac{\lambda_n^2 N^{2(1+\delta)}}{2M^2 n}\right).
\]
Consequently,
\[
\mathbb{P}\bigl(|q_n-U_n|\ge \tfrac{\lambda_n}{\sin\theta}\bigr)
\le
\exp\!\left(-\frac{\lambda_n^2 N^{2(1+\delta)}}{2M^2 n}\right).
\]
\end{lemma}

\begin{prop}\label{prop:randombigO}
Let $\theta=\pi/N$, $x=2\cos\theta$.  
With probability at least $1-\exp(-N^\delta)$,
\begin{align*}
|q_{N-1}-1| &= O\!\left(N^{-\frac12(1+\delta)}\right),\\
|q_{N}| &= O\!\left(N^{-\frac12(1+\delta)}\right),\\
|q_{N+1}+1| &= O\!\left(N^{-\frac12(1+\delta)}\right),\\
|s_{N-2}-2\cos\theta| &= O\!\left(N^{-\frac12(1+\delta)}\right),\\
|s_{N-1}-1| &= O\!\left(N^{-\frac12(1+\delta)}\right).
\end{align*}
\end{prop}

\begin{proof}
Let $\theta=\pi/N$ and $x=2\cos\theta$.  We apply Lemma~\ref{lem:applyAzuma} with the choice
\[
\lambda_n := \sqrt{2n}\,N^{-(3/2+\delta/2)} .
\]
Then for each $1\le n\le N+1$, Lemma~\ref{lem:applyAzuma} gives
\[
\mathbb{P}\bigl(|\delta_n|\ge \lambda_n\bigr)
\le
\exp\!\left(-c\,\frac{\lambda_n^2 N^{2(1+\delta)}}{n}\right)
=
\exp\!\left(-c\,N^{\delta}\right),
\]
for some constant $c>0$ depending only on the uniform bound on $|\xi_k|$.

Using the identity $\sin\theta\,(q_n-U_n)=-\textrm{Im}(\delta_n e^{-in\theta})$, we obtain
\[
|q_n-U_n|\le \frac{|\delta_n|}{\sin\theta}.
\]
Since $\sin(\pi/N)\sim 1/N$, there is a constant $C>0$ such that
\[
|q_n-U_n|\le C N\,|\delta_n|.
\]
Therefore, on the event $\{|\delta_n|\le \lambda_n\}$ we have
\[
|q_n-U_n|\le C N \lambda_n
\le C \sqrt{2n}\,N^{-(1/2+\delta/2)}
=O\!\left(N^{-\frac12(1+\delta)}\right),
\qquad 1\le n\le N+1.
\]

By a union bound over $1\le n\le N+1$,
\[
\mathbb{P}\Bigl(\max_{1\le n\le N+1}|\delta_n|>\lambda_n\Bigr)
\le (N+1)\exp(-cN^\delta)
\le \exp(-c' N^\delta)
\]
for some $c'>0$ and all $N$ large. Hence, with probability at least
$1-\exp(-c' N^\delta)$ we have simultaneously for all $1\le n\le N+1$,
\[
|q_n-U_n| = O\!\left(N^{-\frac12(1+\delta)}\right).
\]

Now use the explicit values of $U_n$ at $\theta=\pi/N$:
\[
U_{N-1}(2\cos(\pi/N))=\frac{\sin((N-1)\pi/N)}{\sin(\pi/N)}=1,\quad
U_N(2\cos(\pi/N))=\frac{\sin(\pi)}{\sin(\pi/N)}=0,
\]
\[
U_{N+1}(2\cos(\pi/N))=\frac{\sin((N+1)\pi/N)}{\sin(\pi/N)}=-1.
\]
This yields the first three estimates:
\[
|q_{N-1}-1|=O\!\left(N^{-\frac12(1+\delta)}\right),\quad
|q_N|=O\!\left(N^{-\frac12(1+\delta)}\right),\quad
|q_{N+1}+1|=O\!\left(N^{-\frac12(1+\delta)}\right).
\]

Finally, the same argument applies to the shifted sequence $s_n$ (as noted in the remark preceding Lemma~\ref{lem:martingale}), giving
\[
|s_n-U_n|=O\!\left(N^{-\frac12(1+\delta)}\right)
\]
uniformly for $n\le N+1$ on the same high-probability event. In particular,
\[
|s_{N-2}-U_{N-2}|=O\!\left(N^{-\frac12(1+\delta)}\right),\qquad
|s_{N-1}-U_{N-1}|=O\!\left(N^{-\frac12(1+\delta)}\right).
\]
Since
\[
U_{N-2}(2\cos(\pi/N))
=\frac{\sin((N-2)\pi/N)}{\sin(\pi/N)}
=\frac{\sin(2\pi/N)}{\sin(\pi/N)}=2\cos(\pi/N),
\]
and $U_{N-1}(2\cos(\pi/N))=1$, we obtain
\[
|s_{N-2}-2\cos(\theta)|=O\!\left(N^{-\frac12(1+\delta)}\right),\qquad
|s_{N-1}-1|=O\!\left(N^{-\frac12(1+\delta)}\right),
\]
as claimed.
\end{proof}

\end{document}
